\centerline{\bf \Huge Обратные остатки и теорема Вильсона.}

\begin{flushright}
        \scriptsize{Занятие: *заканчивается в 20:30*\\
                    20:31 прибыл Герман Сатору}
\end{flushright}

\problem{1}
Решите сравнения (найдите все подходящие $x$ и докажите, что других нет):

а) $5x \equiv 2 \ (mod \ 3)$

б) $3x \equiv 2 \ (mod \ 11 )$

в) $6x \equiv 1 \ (mod \ 13)$


\problem{2}
Какой остаток дает $x + y$ при делении на $17$, если 

a) $x - 16y \equiv 2 \ (mod \ 17);$

б) $3x \equiv 5+14y \ (mod  \ 17)?$

\problem{3}
\textbf{Супер важная задача!} Дано простое число $p$ и его некоторый ненулевой остаток $a$.

a) Докажите, что в последовательности $0 \cdot a, 1 \cdot a, 2 \cdot a, \ldots, (p-1)\cdot a$ все числа дают разные остатки по модулю $p$.

б) Докажите, что существует и при том единственный остаток $b$, что
$ab \equiv 1 \ (mod \ p)$ (такой остаток $b$ называется \textit{обратным} остатка $a$).

в) Какие остатки совпадают со своими обратными остатками?

\problem{4}
a) \textbf{Теорема Вильсона.} Пусть $p$ $-$ некоторое простое число. Докажите, что $(p-1)! \equiv -1 \ (mod \ p)$

\noindent
б) \textbf{Обратная теорема Вильсона.} Докажите, что если $(n-1)! \equiv -1 (mod \ n)$ то число $n \ -$ простое.

\problem{5}
Докажите, что при простом $p$ верно 

\[ \left(\dfrac{p-1}{2}\right)!^2 + (-1)^{\dfrac{p-1}{2}} \equiv -1 \ (mod \ p)\]

\problem{6}
\textbf{Теорема Клемента.} Докажите, что числа $p$ и  $p + 2$  являются простыми числами-близнецами тогда и только тогда, когда  $4((p-1)! + 1) + p \equiv 0 \ (mod \ p^2 + 2p).$